\documentclass[../main.tex]{subfiles}

\begin{document}

\section{第 4 课：数学公式入门}

\subsection*{本课目标}

\begin{itemize}
  \item 区分行内公式和独立公式。
  \item 会写上标、分数、简单函数表达式。
  \item 知道为什么数学模式要用特殊语法。
\end{itemize}

\subsection*{最小可运行示例}

\begin{CodeBlock}
\documentclass[12pt,a4paper]{ctexart}
\usepackage{amsmath} % 常用数学公式包

\begin{document}
行内公式：$a^2 + b^2 = c^2$。

独立公式：
\[
E = mc^2
\]

带编号公式：
\begin{equation}
f(x) = \frac{x^2 + 1}{x + 1}
\end{equation}
\end{document}
\end{CodeBlock}

\subsection*{简明中文解释}

\begin{enumerate}
  \item `\usepackage{amsmath}` 是常用数学扩展包，初学阶段很值得直接带上。
  \item `$ ... $` 表示行内公式，也就是公式嵌在一段话里。
  \item `\[ ... \]` 表示单独成行的公式，默认不编号。
  \item `equation` 环境会给公式自动编号。
  \item `^` 表示上标，`\frac{分子}{分母}` 表示分数。
\end{enumerate}

\subsection*{改什么、看什么}

打开 `\texttt{examples/lesson04\_math.tex}`，先只改公式内容：

\begin{itemize}
  \item 把 `c^2` 改成 `c^3`，观察上标变化。
  \item 把 `\[` 改成普通文本，感受“数学模式”和“普通文字模式”的区别。
\end{itemize}

\subsection*{动手修改任务}

\begin{enumerate}
  \item 把第一个公式改成 `$x_1 + x_2 = 10$`，体验下标写法。
  \item 把独立公式改成你熟悉的一个公式，例如 `y = kx + b`。
  \item 把分数改成 `\frac{1}{2}` 或 `\frac{a+b}{c}`，观察排版。
  \item 在正文里补一句“这个公式表示什么”。
\end{enumerate}

\subsection*{常见报错或易错点}

\begin{itemize}
  \item 少写一个 `$`，很容易导致后面大片内容都进入错误的数学模式。
  \item `\frac` 必须跟两个花括号组，少一个就会报错。
  \item 如果把 `^` 放在普通文本模式里，常会报语法错误或得到意外结果。
\end{itemize}

\subsection*{对应文件}

本课建议直接修改：`\texttt{examples/lesson04\_math.tex}`。
\end{document}
